% p.80 — الباب الرابع والثلاثون في قوانين التعليل وغيره (تابع: القوانين 32–39)
% Layout: twocol. Right column: rules 32–35 and the first line of 36 (ends «ساكنا و اردت»);
% left column: rest of 36 («ان توصله …»), 37–39. Rule 39 breaks off at the foot of the left column
% and continues on p.81.
\mspage{80}
\begin{twocol}
\begin{rightcol}
(32) لا يأتى وزن فَتَحَ يَفْتَحُ الا اذا كان فى العين او اللام حرف من حروف الحلق وهى ستة الهمزة والهاء والحاء والخاء والعين والغين مثل قطع يقطع و طحن يطحن
\par
(33) اذا وقع الواو بين الياء والكسرة اللازمة يحذف للثقل تقول فى يَوْعِدُ يَعِدُ وَحذف فى تَعِدُ اتباعا ليعد وحذف فى يَقَعُ لان اصله يَوْقِعُ ثم فُتِحَ القاف لوجود حرف الحلق
\par
(34) اذا اجتمع حرفان من جنس واحد و كانا متحركين فيدغم للتخفيف تقول فى مَدَدَ مَدَّ ولا يدغم اذا كان الثانى ساكنا مثل يَمْدُدْنَ لان السكون خفيف و لا حاجة للتخفيف
\par
(35) اذا كان حرف العلة متحركا وكان ما قبله ساكنا ينقل حركته الى ما قبله لقوته و لضعف حرف العلة كما مر مثل يَقْوُلُ يَقُوْلُ ويَبْيِعُ يَبِيْعُ ثم يبدل حرف العلة مطابقا لحركة ما قبله للمناسبة فتقول فى يَخْوَفُ يَخَوْفُ ثم يَخَافُ وفى يُعْوِنُ يُعِوْنُ ثم يُعِيْنُ
\par
(36) اذا كان حرف الاخر ساكنا و اردت
\end{rightcol}
\begin{leftcol}
ان توصله الى كلمة ذات الالف واللام فتكسر الحرف الاخر تقول فى احدٌ الله اَحَدُنِ الله وقد يُضم اذا كان ما قبل الآخر مضموما مثل عليكمُ السّلامُ و هُمُ السُّفَهاء هذا فى الضمائر غالب
\par
(37) تكتب الهمزة فى الاول على صورة الالف مثل اَبٍ و اُمٍّ و اِبْلٍ وفى الوسط اذا كانت ساكنة على وفق حركة ما قبلها مثل رَأْسٍ ولُؤْمٍ و ذِئْبٍ واذا كانت متحركة على وفق حركة نفسها مثل سَأَلَ لَؤُمَ سَئِمَ واذا كانت فى اخر الكلمة فعلى وفق حركة ما قبلها لا على حركة نفسها مثل قَرَأَ جَرُؤَ فَتِئَ واذا كان ما قبلها ساكنا فيكتب على الصورة الاصلية مثل خَبْءٍ
\par
(38) الخروج من الكسرة الى الضمة قبيح عندهم ليس عندهم وزن فِعُلُ وكذالك العكس ايضا لكن اجيز فى المجهول لغير العلة للفرق بين المعروف والمجهول نحو نَصَرَ نُصِرَ ولا يجوز فى قَوَلَ قُوِلَ لكونه ثقيلا فى العلة
\par
(39) الواو ضمة تقديرية والياء كسرة
\end{leftcol}
\end{twocol}
